% !TeX root = ../main.tex
\section{Term-by-term correspondence of the single-phase energy balances}
\label{app:single_phase_energy_audit}

This appendix verifies the correspondence between
\eqref{eq:MC_single_phase_energy_limit} and
\eqref{eq:MC_montanaro_electric_enthalpy_limit}.  Take the nonporous,
single-phase specialization used to obtain
\eqref{eq:MC_single_phase_energy_limit}, write
\(\mbf E^{M}=\mbf E\), and identify the phase material derivative with the
single-phase material derivative.  The field kinematics stated before
\eqref{eq:MC_single_phase_energy_limit} give
%
\begin{equation}
\mbf d\cdot\nabla_{\mbf x}\phasedot\varphi
=
-\mbf d\cdot\phasedot{\mbf E}
-
\left(
\mbf E\otimes\mbf d
\right):\mbf L.
\label{eq:appendix_single_phase_field_rate_audit}
\end{equation}
%
The electrical boundary and free-charge terms in
\eqref{eq:MC_single_phase_energy_limit} satisfy
\eqref{eq:MC_electrostatic_power_identity}; hence
%
\begin{equation}
\nabla_{\mbf x}\cdot
\left(
\phasedot\varphi\,\mbf d
\right)
-
\varrho\phasedot\varphi
=
\mbf d\cdot\nabla_{\mbf x}\phasedot\varphi.
\label{eq:appendix_single_phase_charge_power_audit}
\end{equation}
%
Substitution of
\eqref{eq:appendix_single_phase_field_rate_audit} and
\eqref{eq:appendix_single_phase_charge_power_audit} into
\eqref{eq:MC_single_phase_energy_limit}, followed by
\(\bs\sigma_\xi^+=\omega_\xi^+\mbf I+\mbf E\otimes\mbf d\) from
\eqref{eq:vp_electrical_stress_definition}, gives
%
\begin{align}
\rho_\xi\phasedot e_\xi
+\phasedot\omega_\xi^+
&=
\left(
\bs\sigma_\xi'-\bar p_\xi\mbf I
\right):\mbf L_\xi
-\nabla_{\mbf x}\cdot\mbf q_\xi
-\mbf d\cdot\phasedot{\mbf E}
+\rho_\xi r_\xi.
\label{eq:appendix_single_phase_electric_enthalpy_audit}
\end{align}
%
Here the two \(\omega_\xi^+\operatorname{tr}\mbf L_\xi\) terms cancel:
one is part of the material-volume rate on the left of
\eqref{eq:MC_single_phase_energy_limit}, and the other is the isotropic part
of \(\bs\sigma_\xi^+\).  The dyadic Maxwell term cancels the second term on
the right of \eqref{eq:appendix_single_phase_field_rate_audit}.  Thus the
mechanical, electrical, heat-flux, and heat-supply terms have each been
accounted for once.

To compare with the polarization convention of
\cite{montanaro2015thermo_electroelasticity}, separate the vacuum displacement
from the material polarization according to
%
\begin{equation}
\begin{aligned}
\mbf d
&=
\varepsilon_0\mbf E+\mbf P,
&
\bs\pi
&=
\frac{\mbf P}{\rho_\xi},
\\
\rho_\xi
\left(
\varepsilon-\mbf E^{M}\cdot\bs\pi
\right)
&=
\rho_\xi e_\xi+\omega_\xi^+
+\frac{\varepsilon_0}{2}\mbf E\cdot\mbf E,
\\
\bs\tau
&=
\bs\sigma_\xi'-\bar p_\xi\mbf I
+
\left(
\omega_\xi^+
+\frac{\varepsilon_0}{2}\mbf E\cdot\mbf E
\right)\mbf I.
\end{aligned}
\label{eq:appendix_single_phase_energy_partition_audit}
\end{equation}
%
The last two equalities specify the energy and stress identifications between
the two conventions.  The current electric enthalpy contains the vacuum-field
contribution \(-\varepsilon_0\mbf E\cdot\mbf E/2\); adding the vacuum-field
energy converts it to the polarization-based specific energy.  Because
\(\omega_\xi^+\) and the vacuum-field energy are densities per current volume,
material differentiation of the energy identification gives
%
\begin{align}
\rho_\xi
\frac{D_\xi}{Dt}
\left(
\varepsilon-\mbf E^{M}\cdot\bs\pi
\right)
={}&
\rho_\xi\phasedot e_\xi
+\phasedot\omega_\xi^+
+
\left(
\omega_\xi^+
+\frac{\varepsilon_0}{2}\mbf E\cdot\mbf E
\right)
\operatorname{tr}\mbf L_\xi
\notag\\
&+
\varepsilon_0\mbf E\cdot\phasedot{\mbf E}.
\label{eq:appendix_single_phase_energy_partition_rate_audit}
\end{align}
%
Substitution of
\eqref{eq:appendix_single_phase_electric_enthalpy_audit} into
\eqref{eq:appendix_single_phase_energy_partition_rate_audit}, followed by
\(\mbf d=\varepsilon_0\mbf E+\mbf P\) and the stress identification in
\eqref{eq:appendix_single_phase_energy_partition_audit}, yields
%
\begin{align}
\rho_\xi
\frac{D_\xi}{Dt}
\left(
\varepsilon-\mbf E^{M}\cdot\bs\pi
\right)
&=
\bs\tau:\mbf L_\xi
-\nabla_{\mbf x}\cdot\mbf q_\xi
-\mbf P\cdot\phasedot{\mbf E}^{M}
+\rho_\xi r_\xi.
\label{eq:appendix_single_phase_montanaro_recovery}
\end{align}
%
Equation~\eqref{eq:appendix_single_phase_montanaro_recovery} is the
electric-enthalpy form obtained from equations~(10), (12), and (16) of
\cite{montanaro2015thermo_electroelasticity}.  It is
\eqref{eq:MC_montanaro_electric_enthalpy_limit} after the notational
identifications \(\rho=\rho_\xi\), \(\mbf L=\mbf L_\xi\),
\(\mbf q=\mbf q_\xi\), and \(r=r_\xi\), together with the energy and stress
identifications in \eqref{eq:appendix_single_phase_energy_partition_audit}.
The current formulation keeps
the free-charge bulk power and electrical boundary flux visible until
\eqref{eq:MC_electrostatic_power_identity} is applied; the polarization
form instead combines them into the field-rate work.  The two balances are
therefore equivalent under the stated single-phase specialization, vacuum-field
energy partition, and stress identification.
